\documentclass[10pt]{article}
\textheight=9.25in \textwidth=7in \topmargin=-.75in
 \oddsidemargin=-0.25in
\evensidemargin=-0.25in
\usepackage{url}  % The bib file uses this
\usepackage{graphicx} %to import pictures
\usepackage{amsmath, amssymb}
\usepackage{theorem, multicol, color}
\usepackage{gfsartemisia-euler}

\setlength{\intextsep}{5mm} \setlength{\textfloatsep}{5mm}
\setlength{\floatsep}{5mm}
\setlength{\parindent}{0em} % new paragraphs are not indented
\setcounter{MaxMatrixCols}{20}
\usepackage{caption}
\captionsetup[figure]{font=small}


%%%%  SHORTCUT COMMANDS  %%%%
\newcommand{\ds}{\displaystyle}
\newcommand{\Z}{\mathbb{Z}}
\newcommand{\arc}{\rightarrow}
\newcommand{\R}{\mathbb{R}}
\newcommand{\N}{\mathbb{N}}
\newcommand{\Q}{\mathbb{Q}}
\renewcommand{\P}{\mathbb{P}}
\newcommand{\blank}{\underline{\hspace{0.33in}}}
\newcommand{\qand}{\quad and \quad}
\renewcommand{\stirling}[2]{\genfrac{\{}{\}}{0pt}{}{#1}{#2}}
\newcommand{\dydx}{\ds \frac{d y}{d x}}
\newcommand{\ddx}{\ds \frac{d}{d x}}
\newcommand{\dvdx}{\ds \frac{d v}{d x}} 
\newcommand{\solution}{\color{blue}\textit{Solution:} \color{black} }
%%%%  footnote style %%%%

\renewcommand{\thefootnote}{\fnsymbol{footnote}}

\pagestyle{empty}

\begin{document}

\begin{flushright}
Chandler Justice - A02313187
\end{flushright}
\noindent \underline{\hspace{3in}}\\
\textbf{Math 6400} Quiz \#1\\
\textbf{Due:} January 9, 2026 @ 23:59\\

	\noindent Please access the \emph{World Wide Web of EverythingAllTheTime} and apprise yourself of the following categories of Mathematical Philosophies (MPs): Platonism, Logicism, Formalism, and Intuitionism. Also, search for ``Math as a language'' which is arguably a mathematical philosophy whose name essentially describes itself.  
	
	\noindent Please contemplate what implications thinking in the following ways has: (1) Math is a language, (2) Platonism is correct, (3) Logicism is correct, (4) Formalism is correct, (5) Intuitionism is correct.  No need for a response; just ... contemplate.
	
	\begin{enumerate} 
		\item Identify 2 of the 5 MPs that are contrary to one another, meaning you're an idiot if you say ``I'm a bit of an X and a bit of a Y as far as my MP is concerned'' (where X and Y are variables over the set of MPs).\\

        \solution  Formalism and Platonism are at odds with each other in their philosophies of the scope of an
        object. Formalism asserts that \textit{a thing is what it is and nothing more}. For example, a formalist believes
that mathematics is simply the symbols and the rules/axiomatic system we assert on those symbols. This
contrasts directly with Platonism’s stance on objects. A Platonist believes that objects have meaning beyond
their asserted meaning (IE, an axiomatic system has a reality that is beyond our representation and under-
standing). These philosophies embody the \textit{math is invented} and \textit{math is discovered} opinions respectively,
and therefore claiming to possess both philosophies is moronic.
	
		
		\item Do you think Mathematics is invented or discovered?  
        Explain, as best you can.\\

\solution I believe mathematics is invented to formalize systems we observe or imagine existing. For
example, the Greeks utilized measures and lengths to represent units in their mathematical systems and
from there were able to deduce several theorems, postulates, properties, etc regarding geometry which
were applicable to the problems they were trying to solve at that time. This story is largely unchanged
from the way math moves forward today. For example, we often want to model physical phenomena on
computers; however, these phenomena are often defined using continuous differential equations which we
either have no chance of solving in a discrete space like a computer or the amount of compute time required
is prohibitively high. Therefore, we invent new methods for solving these equations in the constraints
inhibited by a computer such as a finite $\{$ difference, element, volume $\}$ method. While finite difference
methods were initially introduced by Euler in the 1700s, they were not popularized until the 1950s when we
started simulating physical phenomena in computers. Work is still done today to improve the fidelity and
accuracy of these methods and to increase the scope of their applications. My motivation for saying all that
is that if we did not have the desire to formalize physical phenomena using differential equations and
then further simulating those phenomena in a computer, these mathematical techniques and formalizations
may not have come about; therefore, we can see that these constructs are invented and are pure thought
stuff.\\

Beyond this, mathematics often introduces ideas/concepts where applications are not known upon inception, but these ideas are built on the foundation of existing mathematics and are invented as extensions of existing formalisms by the mathematician. 
	
		\item Which number has a more well-defined existence: $1$ or $\sqrt{2}$.
		Please justify your response.
		
\solution The set of integers and by extension the natural numbers, in my opinion, is a set which better
reflects reality than the complex or real numbers. With this assertion, the number $1$ naturally has a more
well defined existence as it is a member of the natural numbers while $\sqrt{2}$ is not. With that said, we can
observe $\sqrt{2}$ in the real world in places such as a right triangle with $a = b = 1$, meaning $c = \sqrt{2} \approx 1.41$.
Even if we believe $\sqrt{2}$ is a ``weaker'' number by virtue of being irrational, there are many places where it
is an absolutely necessary representation such as in geometric and algebraic settings. With that said, by
virtue of $1$ being a discrete object, it is able to serve more roles as a numeric object versus $\sqrt{2}$. For example,
you cannot have $\sqrt{2}$ slices of pizza or $\sqrt{2}$ of anything for that matter.



		\item A {\bf prime number} is a positive integer greater than $1$ that has only itself and $1$ as divisors\footnote{Suppose $a$ and $b$ are integers; $a$ is a {\bf divisor} of $b$ if $b = ak$, for some integer $k$.}.  
		Is the following argument (think of it as a proof) correct?  If so, what has it proved?  
		Please justify your response.
		
		\emph{Maybe a Proof of Something:} Suppose $A,B$ and $C$ are all the prime numbers (and that they're all different).  
		Compute $N = ABC+1$. Either $N$ is prime or it's not.  
		If it's prime, then this contradicts the assumption that $A,B,$ and $C$ are all the prime numbers. 
		If it is not prime, then it has $A,B,$ or $C$ as a divisor. 
		But the reader can verify that neither $A$ nor $B$ nor $C$ can be a divisor of $N$.
		Therefore, $A,B,$ and $C$ do not comprise the entirety of prime numbers; in particular, there is another prime number, say $D$. 
		The argument may be repeated to conclude $A,B,C,$ and $D$ are not all the prime numbers.\\
		
        \solution The intent with this proof is to show there are always more primes than a set of primes contains;
however, it is not motivated why if the number $N$ is not prime, then it must have $A, B$ or $C$ as a divisor.
This proof is almost correct; however, the phrase ``If it is not prime, then it has $A, B$ or $C$ as a divisor''
should be replaced with something like ``$N$’s prime factorization will contain a prime $D$ which is unique
from $A, B$ and $C$''. With this fix, this proof is now suitable for showing that given a set of prime numbers,
there is always another prime number not included in the provided set.
		
		\item  Without using any technology (other than something to write with), compute, with as much precision as you can, the number $x$ such that $x^2 = 2$.   
		
            \solution From algebra classes, we know $(\sqrt{2})^2 = 2$; however, we want to determine the actual numerical
            value of $\sqrt{2}$ with the tightest bounds we can set (as was explained in class). We can start with initial
            bounds $1 < \sqrt{2} < 2$ since $1^2 = 1$ and $2^2 = 4$. We can tighten our bounds by lowering the upper bound and
            increasing the lower bound until we find the tightest bound we can bother to compute by hand\footnote{in my last submission, I did it by hand. With the new information I can use a calculator, I will try to be more precise in this submission}. I know $\sqrt{2} \approx 1.41$ from memory, so we can tighten from there. We can see $(1.41)^2 = 1.9881$ and $(1.42)^2 = 2.0164$,
            so $1.41 < \sqrt{2} < 1.42$. We can see $(1.417)^2 = 2.007889$ and $(1.412)^2 = 1.993744$. Let's zero in on the last decimal place and determine where the line is between a value starting with $1.9$ and a value starting with $2.0$. Using a calculator now, we can see $(1.414)^2 = 1.999396$ and $(1.415)^2 = 2.002225$ meaning $1.414 < \sqrt{2} < 1.415$. We should add additional precision to our lower bound to get a closer approximation, and then replace our upper bound with this new approximation. We then find $(1.4143)^2 = 2.00024449$, so $1.4140 < \sqrt{2} < 1.4143$. After a few additional iterations, we find the approximation $(1.414214)^2 = 2.000001238$. If we subtract $0.000001$ from this upper bound, we get $(1.414213)^2 = 1.999998409$ as a new lower bound. This gives us the final approximation I am going to compute of $1.414213 < \sqrt{2} < 1.414214$.\\

            This is more than enough precision for pretty much any application of $\sqrt{2}$. We can compute our error using $\frac{b - a}{2}$ where $b$ is the upper bound and $a$ is the lower bound to be $1.4142135 \times 10^{-6}$. This is extremely low error and therefore our bounds serve as a suitable approximation of $\sqrt{2}$.
		
	\end{enumerate}
	


\noindent \underline{\hspace{3in}}\\

\end{document}

